\documentclass{packages/icmc}
\usepackage[utf8]{inputenc}
\usepackage[T1]{fontenc}
\usepackage{geometry}
\usepackage{graphicx}
\usepackage{float}
\usepackage{hyperref}
\usepackage[english,indonesian]{babel}

\title{Optimasi Pemrosesan Template LaTeX Menggunakan Pendekatan Berbasis AI}
\author{Academic Support Agent}
\date{\today}

\begin{document}
\maketitle

\begin{abstract}

\end{abstract}

\chapter{PENDAHULUAN}


\chapter{TINJAUAN PUSTAKA}


% Original content (commented out):
% % ------------------------------------------------------------------------
% ------------------------------------------------------------------------
% Modelo ICMS adaptado por Elias J.
% Mestrado UFMG
% ------------------------------------------------------------------------
% ------------------------------------------------------------------------

% Opções: 
%   Qualificação         = qualificacao 
%   Curso                = doutorado/mestrado
%   Situação do trabalho = pre-defesa/pos-defesa (exceto para qualificação)
% -- opções do pacote babel --
% Idioma padrão = brazil
	%spanish,			% idioma adicional para hifenização
	%english,			% idioma adicional para hifenização
	%brazil				% o último idioma é o principal do documento
\documentclass[mestrado, pre-defesa, spanish, english, brazil]{packages/icmc}

\usepackage{chngcntr}
\usepackage{fontawesome}
\counterwithin{table}{chapter}
%\counterwithin{table}{section}
\counterwithin{figure}{chapter}
%\counterwithin{algorithm}{chapter}
\counterwithin{equation}{chapter}
\counterwithout{footnote}{chapter}


\usepackage{nomencl}
%\usepackage{fontspec}

%\usepackage{fontspec} % already loaded by unicode-math
    

% ---
% Pacotes Opcionais
% ---


\usepackage{xspace} 
\usepackage{rotating}           % Usado para rotacionar o texto
\usepackage[all,knot,arc,import,poly]{xy}   % Pacote para desenhos gráficos
% Este pacote pode conflitar com outros pacotes gráficos como o ``pictex''
% Então é necessário usar apenas um dos pacotes conflitantes

\usepackage[T1]{fontenc}		% Selecao de codigos de fonte.

\usepackage{pdfpages} %insert page

\usepackage{epigraph}

\usepackage{color}				% Controle das cores
\usepackage{subfig}
\usepackage{graphicx}			% Inclusão de gráficos
\usepackage{float}
%\usepackage[bottom]{footmisc} % resolve o problema quando tem uma figura próximo a uma nota de roda pé

\usepackage{media9} %inclusao de gifs
%\usepackage{geometry}
%\usepackage{animate}

\graphicspath{{images/}}


\usepackage{cite}
\usepackage{upquote}

\newcommand{\VerbL}{0.52\textwidth}
\newcommand{\LatL}{0.42\textwidth}

\usepackage{xcolor}
\usepackage{tikz}
\usepackage{setspace} % Line Spacing

%para usar os algoritmos e matematica
\newcommand{\var}{\texttt}

\usepackage{amssymb}
\usepackage{amsmath}
\usepackage{lastpage}


\renewcommand{\citeauthor}[1]{\citeauthoronline{#1} [\citeyear{#1}]} \def\citep{\cite} 
\newcommand{\citeyearpar}[1]{(\citeyear{#1})} \ifx\AbntCitetype\AbntCitetypeALF \def\citet{\citeonline}

%texto
\usepackage[brazil]{babel}
\usepackage{blindtext}

\usepackage{etoolbox}


%\usepackage{siunitx} %para números com notação científica x10-3


%%%%%%%%%%%%%%%%%%%%%%%%%%%%%%%%%%%
% Algoritmos
% \usepackage[Algoritmo]{algorithm}
% %\usepackage[end]{algpseudocode}
% %\usepackage[portuguese,ruled,linesnumbered,lined]{algorithm2e}	%Permite criar algoritmos com numeração
% %\usepackage{algcompatible}
%\usepackage{algorithmicx}
%\usepackage{algorithm}
\usepackage[noend]{algpseudocode}
\usepackage{algorithm2e}

\SetKwInput{Dados}{Dados}
\SetKwInput{Entrada}{Entrada}
\SetKw{Retorna}{retorna}
\SetKwInput{Resultado}{Resultado}
\SetKwInput{Resultado}{Resultado}
\SetKwFor{Para}{para}{faça}{fim para}
\SetKwFor{Enquanto}{Enquanto}{faça}{fim-enquanto}
\SetKwFor{FacaParalelo}{Faça em paralelo enquanto $t \leq t_s$:}{}{fim-paralelo}
 \SetKwIF{Se}{SenaoSe}{Senao}{Se}{então}{Senão se}{Senão}{fim se}
\SetKwFor{ParaTodo}{para todo}{faça}{fim para todo}
\SetKwFor{ParaCada}{para cada}{faça}{fim para cada}
\SetKwProg{Fn}{Função}{}{}
 
\newcommand\mycommfont[1]{\footnotesize\ttfamily\textcolor{blue}{#1}}
\SetCommentSty{mycommfont}

%\usepackage[portuguese]{algorithm2e}

\algblockdefx{FORALLP}{ENDFAP}[1]%
  {\textbf{faça em paralelo:}}%
  {\textbf{fim paralelo}}
% declaration of the new block
\algblock{FORALLP}{ENDFAP}

\algblockdefx{WHILEp}{ENDWHILEp}[1]%
  {\textbf{enquanto} VERDADE \textbf{faça:}}%
  {\textbf{fim}}
% declaration of the new block
\algblock{WHILEp}{ENDWHILEp}

\algrenewcommand\algorithmicwhile{\textbf{enquanto}}
\algrenewcommand\algorithmicdo{\textbf{faça}}
\algrenewcommand\algorithmicend{\textbf{fim}}
\algrenewtext{EndWhile}{\algorithmicend\ \algorithmicwhile}
%%%%%%%%%%%%%%%%%%%%%%%%%%%%%%%%%%%

%\usepackage[figurewithin=none]{caption}
\usepackage{caption}
\usepackage{booktabs}
\usepackage{array}% http://ctan.org/pkg/array
\usepackage{tabularx}
\usepackage{adjustbox} %ajusta tabela
\usepackage{multirow} % para colocar multiplas colunas e linhas
\newcolumntype{P}[1]{>{\centering\arraybackslash}p{#1}}
\newcolumntype{M}[1]{>{\centering\arraybackslash}m{#1}}
\newcolumntype{V}{>{\centering\arraybackslash} m{.1\linewidth} }
%\newcolumntype{M}{>{\centering\arraybackslash}m{\dimexpr.25\linewidth-2\tabcolsep}}

\usepackage{array} % for defining a new column type
\usepackage{varwidth} %for the varwidth minipage environment

% ---
% Informações de dados para CAPA e FOLHA DE ROSTO
% ---
% Tanto na capa quanto nas folhas de rosto apenas a primeira letra da primeira palavra (ou nomes próprios) devem estar em letra maiúscula, todas as demais devem ser em letra minúscula.
\tituloPT{Meu título}
\tituloEN{Title}
\autor[Autor abreviado]{Autor completo}
\genero{M} % Gênero do autor (M = Masculino / F = Feminino)
\orientador[Orientador]{Prof. Dr.}{Guilherme Augusto Silva Pereira}
%\coorientador{Prof. Dr.}{Fulano de Tal}
\curso{PPGEE}
\data{20}{10}{2017} % Data do depósito

% ---


% ---
% RESUMOS
% ---


% Resumo em PORTUGUÊS
% conter no máximo 500 palavras
% conter no mínimo 1 e no máximo 5 palavras-chave
\textoresumo[brazil]{
    meu resumo
    }{robôs }


% resumo em INGLÊS
% conter no máximo 500 palavras
% conter no mínimo 1 e no máximo 5 palavras-chave
\textoresumo[english]{
    my abstract
    }{ robots}
    
% ---
% Configurações de aparência do PDF final
% ---
% alterando o aspecto da cor azul
\definecolor{blue}{RGB}{41,5,195}

% informações do PDF
\makeatletter
\hypersetup{
     	%pagebackref=true,
		pdftitle={\@title}, 
		pdfauthor={\@author},
    	pdfsubject={\imprimirpreambulo},
	    pdfcreator={LaTeX with abnTeX2},
		pdfkeywords={\palavraschave}, 
		colorlinks=true,       		% false: boxed links; true: colored links
    	linkcolor=blue,          	% color of internal links
    	citecolor=blue,        		% color of links to bibliography
    	filecolor=magenta,      	% color of file links
		urlcolor=blue,
		bookmarksdepth=4
}

% Fixing the appearance of the footer in the empty pages in between chapters.
\let\origdoublepage\cleardoublepage
\renewcommand{\cleardoublepage}{%
  \clearpage
  {\pagestyle{empty}\origdoublepage}%
}
% Finished fixing appearance

\makeatother
% --- 

%linhas juntas
%\brokenpenalty=10000
%\clubpenalty=100000 % Para evitar linhas orfas
%\widowpenalty=5 % Para evitar linhas viuvas

%\hyphenpenalty=10000 % Para não hifenizar


% ----------------------------------------------------------
% ELEMENTOS PRÉ-TEXTUAIS
% ----------------------------------------------------------

% Inserir a ficha catalográfica
%USAR AQUI
%acrescentar pdf fornecido pela biblioteca
%\incluifichacatalografica*{tex/pre-textual/ficha_catalografica.pdf}



\incluifichacatalografica{866} % Código Cutter: número atribuído ao sobrenome do autor. Para obtê-lo, consulte a tabela Cutter Sanborn (em http://www.davignon.qc.ca/cutter1.html), procure pelo sobrenome ou forma mais próxima ao sobrenome completo e coloque o número indicado como parâmetro.


% DEDICATÓRIA / AGRADECIMENTO / EPÍGRAFE
%\textodedicatoria*{tex/pre-textual/dedicatoria}
\textoagradecimentos*{tex/pre-textual/agradecimentos}
\textoepigrafe*{tex/pre-textual/epigrafe}

% Inclui a lista de figuras
\incluilistadefiguras

% Inclui a lista de tabelas
\incluilistadetabelas

% Inclui a lista de quadros
%\incluilistadequadros

% Inclui a lista de algoritmos
\incluilistadealgoritmos

% Inclui a lista de códigos
%\incluilistadecodigos

% Inclui a lista de siglas e abreviaturas
\incluilistadesiglas

% Inclui a lista de símbolos
%\incluilistadesimbolos

% ----
% Início do documento
% ----
\begin{document}
% ----------------------------------------------------------
% ELEMENTOS TEXTUAIS
% ----------------------------------------------------------
\textual


\pagenumbering{arabic}
\setcounter{page}{16}
\chapter{Introdução}
\label{chapter:introducao}
%\epigraph{Não importa quão limitado possa parecer o começo: aquilo que é feito uma vez está feito para sempre. }{Henry David Thoreau, 1817-1862}
\input{tex/1_introducao.tex}


\chapter{Revisão bibliográfica}
\label{chapter:revisao}
\input{tex/2_revisaoBibliografica.tex}



\chapter{Metodologia}
\label{chapter:metodologia}
\input{tex/3_metodologia.tex}


\chapter{Resultados}
\label{chapter:resultados}
\input{tex/4_resultados.tex}

\chapter{Conclusões e trabalhos futuros}
\label{chapter:conclusao}
\input{tex/5_conclusao.tex}





% ----------------------------------------------------------
% ELEMENTOS PÓS-TEXTUAIS
% ----------------------------------------------------------
\postextual

% ----------------------------------------------------------
% Referências bibliográficas
% ----------------------------------------------------------
%\bibliographystyle{amsplain}
\bibliography{references}



% ---------------------------------------------------------------------
% GLOSSÁRIO
% ---------------------------------------------------------------------

% Arquivo que contém as definições que vão aparecer no glossário
%\input{tex/glossario}
% Comando para incluir todas as definições do arquivo glossario.tex
%\glsaddall
% Impressão do glossário
%\printglossaries

% ----------------------------------------------------------
% Apêndices
% ----------------------------------------------------------

% ---
% Inicia os apêndices
% ---
\begin{apendicesenv}

    \chapter{XXXXX }
    \label{chapter:diagramas}
    \input{tex/appendix/diagramas}
    
    % \chapter{Configuração do programa JabRef}
    % \label{chapter:configuracao-jabref}
    % \input{tex/appendix/configuracao-jabref}

\end{apendicesenv}
% ---


% ----------------------------------------------------------
% Anexos
% ----------------------------------------------------------

% ---
% Inicia os anexos
% ---
%\begin{anexosenv}
%
%    \chapter{Forma rápida de interpretar o diagrama de caixa} 
%    \label{anexo:diagrama}
%    \input{tex/annex/diagrama}
%
%\end{anexosenv}
% ---

\end{document}

% ---
% Finaliza a parte no bookmark do PDF, para que se inicie o bookmark na raiz
% ---
\bookmarksetup{startatroot}% 
% ---

\end{document}